Krátké shrnutí a cíl sekce: Popis praktických nástrojů a pracovních postupů pro generování a vyhodnocení cvičných úloh v matematice s důrazem na sběr postupu řešení a diagnostiku chyb pomocí nástrojů z rodiny Microsoft (Pokrok v matematice, OneNote, math.microsoft.com) \cite{kb_ai_101_tipu_pdf_04e4a115_page_8_1}.

Praktický pracovní postup pro jeden cyklus domácího úkolu
\begin{enumerate}
  \item Výběr oblasti a generování úloh: učitel v rozhraní „Pokrok v matematice“ vybere oblast, pro kterou chce vygenerovat příklady; systém nabídne sadu úloh, ze které učitel vybírá \cite{kb_ai_101_tipu_pdf_04e4a115_page_8_1}.
  \item Konfigurace požadavků na odevzdání: při nastavování zadání lze určit, zda studenti musí nahrát i postup řešení a zda mají studenti hodnotit obtížnost příkladů \cite{kb_ai_101_tipu_pdf_04e4a115_page_8_1}.
  \item Odevzdání a předběžné zpracování AI: studenti vypracují úlohy a přiloží postupy; AI nástroj provede předopravování (označí pravděpodobné chyby, doplní oblast, kde došlo k chybě) a připraví přehled pro učitele \cite{kb_ai_101_tipu_pdf_04e4a115_page_8_1}.
  \item Pedagogická revize a diagnostika: učitel rychle projde označené případy, upraví hodnocení a použije výsledky k cílenému doplnění výuky (skupinová remediace, individuální úkoly) \cite{kb_ai_101_tipu_pdf_04e4a115_page_8_1}.
\end{enumerate}

Hlavní nástroje a funkcionality (rychlý přehled)
\begin{itemize}
  \item Pokrok v matematice — generování cvičných sad, možnost vyžadovat nahrání postupu, předopravování a přehledy o silných/slabých stránkách žáků \cite{kb_ai_101_tipu_pdf_04e4a115_page_8_1}.
  \item OneNote — převod rukopisných rovnic do digitální podoby, včetně zobrazení postupu a podpory výpočtu či grafu \cite{kb_ai_101_tipu_pdf_04e4a115_page_8_1}.
  \item math.microsoft.com — webové rozhraní slučující nástroje pro práci s matematickými úlohami a cvičeními \cite{kb_ai_101_tipu_pdf_04e4a115_page_8_1}.
  \item Rodina Microsoft školních aplikací — tyto funkce jsou součástí ekosystému školních licencí (např. Office 365 / školní varianty), ověřte dostupnost pro vaši instituci \cite{kb_ai_101_tipu_pdf_04e4a115_page_3_1}.
\end{itemize}

Návrhy pro zařazení do domácích úkolů a diagnostiky (praktické tipy)
\begin{itemize}
  \item Vyžadujte nahrání postupu řešení tam, kde je cílem hodnotit postupy a porozumění; v takových úlohách je „postup důležitější než výsledek“ — možnost požadovat postup lze nastavit v Pokrok v matematice \cite{kb_ai_101_tipu_pdf_04e4a115_page_8_1}.
  \item Využijte automatické označení chyb AI jako predikci oblastí, kde studenti potřebují podporu; plánujte krátké remediace založené na těchto datech \cite{kb_ai_101_tipu_pdf_04e4a115_page_8_1}.
  \item Kombinujte digitální převod rukopisu (OneNote) s požadavkem na nahrání fáze řešení — usnadní to automatizované zpracování a archivaci postupů \cite{kb_ai_101_tipu_pdf_04e4a115_page_8_1}.
  \item Před nasazením zkontrolujte dostupnost funkcí ve školních licencích (Office 365 A1/A3/A5 apod.) a informujte studenty o pravidlech zpracování odevzdaných dat \cite{kb_ai_101_tipu_pdf_04e4a115_page_3_1}. (general background knowledge: konkrétní nastavení retence a přístupů nastaví správa IT fakulty.)
\end{itemize}

Krátký checklist před prvním použitím
\begin{itemize}
  \item Ověřit dostupnost Pokrok v matematice, OneNote a příslušných funkcí v používané univerzitní licenci \cite{kb_ai_101_tipu_pdf_04e4a115_page_3_1}.
  \item Rozhodnout, kdy požadovat nahrání postupu (procesní úlohy) a jasně to specifikovat v instrukcích pro studenty \cite{kb_ai_101_tipu_pdf_04e4a115_page_8_1}.
  \item Nastavit, kdo má přístup k datům o výkonu studentů a jak budou použita pro diagnostiku (general background knowledge).
  \item Otestovat celý tok (vygenerovat úlohy, odevzdat testovací řešení, projít AI označeními) před nasazením do skutečné třídy \cite{kb_ai_101_tipu_pdf_04e4a115_page_8_1}.
\end{itemize}

Poznámka na závěr: Popisované postupy a nástroje vycházejí z implementací a funkcí rodiny vzdělávacích nástrojů Microsoft (Pokrok v matematice, OneNote, math.microsoft.com) uvedených v dostupné literatuře \cite{kb_ai_101_tipu_pdf_04e4a115_page_8_1,kb_ai_101_tipu_pdf_04e4a115_page_3_1}. Pro konkrétní technickou integraci a nastavení ochrany dat spolupracujte s IT oddělením fakulty (general background knowledge).